\subsection{PLONK, a zkSNARK protocol}
\label{plonk}

\subsubsection{Arithmetization}
\label{arith-intro}

The promise of programmable cryptography is that we should be able to
perform proofs for arbitrary functions. That means we need a
``programming language'' that we will write our function in. For PLONK,
the choice that is used is: \textbf{systems of quadratic equations over
\({\mathbb{F}}_{q}\)}. In other words, PLONK is going to give us the
ability to prove that we have solutions to a system of quadratic
equations.

strong(Situation):\\

This leads to the natural question of how a function like SHA-256 can be
encoded into a system of quadratic equations. This process of encoding a
problem into algebra is called \emph{arithmetization}. It turns out that
quadratic equations over \({\mathbb{F}}_{q}\), viewed as an NP problem
called Quad-SAT, is \emph{NP-complete}; in other words, any NP problem
can be rewritten as a system of quadratic equations. If you are not
familiar with this concept, the upshot is that Quad-SAT being
NP-complete means it can serve as a reasonable arithmetization that can
express most reasonable (NP) problems.

strong(Remark):\\
\phantomsection\label{np-complete}{}

So for example, any NP decision problem should be encodable. Still, such
a theoretical reduction might not be usable in practice: Polynomial
factors might not matter in complexity theory, but they do matter a lot
to engineers and end users.

But it turns out that Quad-SAT is actually reasonably codeable. This is
the goal of projects like Circom\footnote{https://docs.circom.io/},
which gives a high-level language that compiles a function like SHA-256
into a system of equations over \({\mathbb{F}}_{q}\) that can be used in
practice. Systems like this are called \emph{arithmetic circuits}, and
Circom is appropriately short for ``circuit compiler''. If you are
curious, you can see how SHA-256 is implemented in Circom on
GitHub.\footnote{https://github.com/iden3/circomlib/blob/master/circuits/sha256/sha256.circom}

So, the first step in proving a claim like ``I have a message \(M\) such
that \(\operatorname{sha}(M) = \ 0xa91af3ac...\)'' is to translate the
claim into a system of quadratic equations. This process is called
``arithmetization.''

One approach (suggested by \hyperref[np-complete]{{[}np-complete{]}}) is
to represent each bit involved in the calculation by a variable
\(x_{i}\) (which would then be constrained to be either 0 or 1 by an
equation \(x_{i}^{2} = x_{i}\)). In this setup, the value \(0xa91af3ac\)
would be represented by 32 public bits \(x_{1},\ldots,x_{32}\); the
unknown message \(M\) would be represented by some private variables;
and the calculation of \(\operatorname{sha}\) would introduce a series
of constraints, maybe involving some additional variables.

We will not get into any more details of arithmetization here.

\subsubsection{An instance of PLONK}

PLONK is going to prove solutions to systems of quadratic equations of a
very particular form:

strong(Definition):\\

strong(Remark):\\

Back to PLONK. Our protocol needs to do the following: Peggy and Victor
have a PLONK instance given to them. Peggy has a solution to the system
of equations, i.e., an assignment of values to each \(a_{i}\),
\(b_{i}\), \(c_{i}\) such that all the gate constraints and all the copy
constraints are satisfied. Peggy wants to prove this to Victor
succinctly and without revealing the solution itself. The protocol then
proceeds by having:

\begin{enumerate}
\item
  Peggy sends a polynomial commitment corresponding to \(a_{i}\),
  \(b_{i}\), and \(c_{i}\) (the details of what polynomial are described
  below).
\item
  Peggy proves to Victor that the commitment from Step 1 satisfies the
  gate constraints.
\item
  Peggy proves to Victor that the commitment from Step 1 also satisfies
  the copy constraints.
\end{enumerate}

Let us now explain how each step works.

\subsubsection{Step 1: the commitment}

In PLONK, we will assume that \(q \equiv 1\operatorname{mod}n\), which
means that we can fix \(\omega \in {\mathbb{F}}_{q}\) to be a primitive
\(n\)-th root of unity.

Then, by polynomial interpolation, Peggy chooses polynomials \(A(X)\),
\(B(X)\), and \(C(X)\) in \({\mathbb{F}}_{q}\lbrack X\rbrack\) such that
\[A\left( \omega^{i} \right) = a_{i},B\left( \omega^{i} \right) = b_{i},C\left( \omega^{i} \right) = c_{i}\text{  for all  }i = 1,2,\ldots,n.\]
\phantomsection\label{plonk-setup}{}

We specifically choose \(\omega^{i}\) because that way, if we use
{[}root-check{]} on the set
\(\left\{ \omega,\omega^{2},\ldots,\omega^{n} \right\}\), then the
polynomial called \(Z\) is just
\(Z(X) = (X - \omega)\ldots\left( X - \omega^{n} \right) = X^{n} - 1\),
which is really nice. In fact, often \(n\) is chosen to be a power of
\(2\) so that \(A\), \(B\), and \(C\) are very easy to compute, using a
fast Fourier transform. (Note: When you are working in a finite field,
the fast Fourier transform is sometimes called the ``number theoretic
transform'' (NTT) even though it is exactly the same as the usual FFT.)

Then:

strong(Algorithm):\\
``Commitment step of PLONK'':

\begin{enumerate}
\item
  Peggy interpolates \(A\), \(B\), \(C\) as in
  \hyperref[plonk-setup]{{[}plonk-setup{]}}.
\item
  Peggy sends \(\operatorname{Com}(A)\), \(\operatorname{Com}(B)\),
  \(\operatorname{Com}(C)\) to Victor.
\end{enumerate}

To reiterate, each commitment is a single value -- a 256-bit elliptic
curve point -- that can later be ``opened'' at any value
\(x \in {\mathbb{F}}_{q}\).

\subsubsection{Step 2: gate check}

Both Peggy and Victor know the PLONK instance, so they can interpolate a
polynomial \(Q_{L}(X) \in {\mathbb{F}}_{q}\lbrack X\rbrack\) of degree
\(n - 1\) such that
\[Q_{L}\left( \omega^{i} \right) = q_{L,i}\quad\text{ for  }i = 1,\ldots,n.\]
The analogous polynomials \(Q_{R}\), \(Q_{O}\), \(Q_{M}\), \(Q_{C}\) are
defined in the same way.

Now, what do the gate constraints amount to? Peggy is trying to convince
Victor that the equation \[\begin{aligned}
Q_{L}(x)A(x) + Q_{R}(x)B(x) + Q_{O}(x)C(x) \\
\phantom{0} + Q_{M}(x)A(x)B(x) + Q_{C}(x) & = 0
\end{aligned}\] \phantomsection\label{plonk-gate}{}

is true for the \(n\) numbers
\(x = \omega,\omega^{2},\ldots,\omega^{n}\).

However, Peggy has committed \(A\), \(B\), \(C\) already, while all the
\(Q_{\ast}\) polynomials are globally known. So this is a direct
application of {[}root-check{]}:

strong(Algorithm):\\
Gate check:

\begin{enumerate}
\item
  Both parties interpolate five polynomials
  \(Q_{\ast} \in {\mathbb{F}}_{q}\lbrack X\rbrack\) from the \(5n\)
  coefficients \(q_{\ast}\) (globally known from the PLONK instance).
\item
  Peggy uses {[}root-check{]} to convince Victor that
  \hyperref[plonk-gate]{{[}plonk-gate{]}} holds for \(X = \omega^{i}\)
  (that is, the left-hand side is indeed divisible by
  \(Z(X) := X^{n} - 1\)).
\end{enumerate}

\subsubsection{Step 3: proving the copy
constraints}\label{copy-constraint-deferred}

The copy constraints are the trickiest step. There are a few moving
parts to this idea, so we skip it for now and dedicate
{[}copy-constraints{]} to it.

\subsubsection{Step 4: public and private witnesses}

The last thing to be done is to reveal the value of public witnesses, so
the prover can convince the verifier that those values are correct. This
is simply an application of {[}root-check{]}. Let us say the public
witnesses are the values \(a_{i}\), for all \(i\) in some set \(S\). (If
some of the \(b\)`s and \(c\)`s are also public, we will just do the
same thing for them.) The prover can interpolate another polynomial,
\(A^{\text{public}}\), such that
\(A^{\text{public }}\left( \omega^{i} \right) = a_{i}\) if \(i \in S\),
and \(A^{\text{public }}\left( \omega^{i} \right) = 0\) if
\(i \notin S\). Actually, both the prover and the verifier can compute
\(A^{\text{public}}\), since all the values \(a_{i}\) are publicly
known!

Now the prover runs {[}root-check{]} to prove that
\(A - A^{\text{public}}\) vanishes on \(S\). (And similarly for \(B\)
and \(C\), if needed.) And we are done.
